\begin{theorem}[Conjecture of Papadakis and Petrotou] \label{thm-PPconj}
Let $\Delta$ be a  pseudo-manifold of dimension $n-1$ over $K$. For any integer vectors $I, J$ with $n$ components
\[ \partial_I W_\Delta(x_J) = \begin{cases} 
(W_\Delta (\sqrt{x_I x_J}))^2 & \text{if $\sqrt{x_I x_J}$ exists,} \\
0 & \text{otherwise.}
\end{cases}
\]
\end{theorem}

We will prove the conjecture below. Let us first see that it implies Theorem~\ref{thm-lef}. The argument here is similar to the proof of Theorem~\ref{thm-PP} in \cite{PapadakisPetrotou}. In fact, it is very natural to extend this theorem to the case of pseudo-manifolds $\Delta$ as in \cite{APP}. Recall that in Section~\ref{sec-mixed} we defined for any oriented pseudo-manifold $\Delta$ the algebra $\bar{H}(\Delta)$. This is equal to the algebra $H(\Delta)$ if $\Delta$ is a homology sphere.

\begin{coro} \label{cor-PPconj}
Let $h \in K[x_1,\ldots, x_N]_m$ for  some $0 \leq m \leq n/2$, and let $I, J$ have $n, n - 2m$ components, respectively. Then
\[ \partial_I W_\Delta(h^2 x_J) = (W_\Delta(h \sqrt{x_I x_J}))^2 \]
if $\sqrt{x_I x_J}$ exists, and is otherwise zero.
\end{coro}

\begin{proof}
Write $h$ as a linear combination of monomials, $h=\sum_L b_L x_L$. Then
\[ W_\Delta(h^2 x_J) = W_\Delta ( \sum_L b_L^2 x_L^2 x_J) = \sum_L b_L^2 W_\Delta (x_L^2 x_J).\]
Applying the derivative $\partial_I$ to this and using  Theorem~\ref{thm-PPconj}, we get
\[  \sum_L b_L^2 \partial_I W_\Delta (x_L^2 x_J) = \sum_L b_L^2 (W_\Delta (x_L \sqrt{x_I x_J}))^2 = (W_\Delta ( \sum_L b_L x_L \sqrt{x_I x_J}))^2\]
if the square root exists, and zero otherwise.
\end{proof}

The previous corollary shows why Theorem~\ref{thm-PPconj} is well suited for proving anisotropy theorems in characteristic $2$. The expression $W_\Delta(h \sqrt{x_I x_J})$ on the right hand side is the \Po pairing between $h$ and $\sqrt{x_I x_J}$. If $W_\Delta(h^2 x_J)$ on the left hand side is zero, then $h$ is orthogonal to $\sqrt{x_I x_J}$ for all $I$.

\begin{coro} \label{cor-reduce}
Let $\Delta$ be a pseudo-manifold of dimension $n-1$ over $K$. Consider $g\in K[x_1,\ldots, x_N]_m$ such that 
\[ l^p g^2 = 0 \quad \text{in $\bar{H}^{p+2m}(\Delta)$}\]
for some  $0\leq p \leq n-2m$. 
Let $q = \lfloor \frac{p}{2}\rfloor$. Then
\[ l^q g = 0 \quad \text{in $\bar{H}^{q+m}(\Delta)$.}\]
\end{coro}

\begin{proof}
If $l^p g^2 = 0$ in $\bar{H}^{p+2m}(\Delta)$, then for any integer vectors $I$ and $J$ with respectively $n$ and $n-2m-p$ components,
\[ \partial_I W_\Delta(l^p g^2 x_J) = 0. \]
If $p$ is even, then $p/2 = q$ and by Corollary~\ref{cor-PPconj},
\[ \partial_I W_\Delta(l^p g^2 x_J) = \partial_I W_\Delta((l^q g)^2 x_J) = (W_\Delta( l^q g \sqrt{x_I x_J}))^2 = 0 \]
if the square root exists, and so
\[ W_\Delta( l^q g \sqrt{x_I x_J}) = 0. \]
Since $I$ and $J$ can be chosen such that $x_I x_J = x_L^2$  for any of the monomials $x_L$ generating $\bar{H}^{n-q-m}(\Delta)$,  we conclude that $l^q g = 0$ in $\bar{H}^{q+m}(\Delta)$.

If $p$ is odd, then $l^p g^2 = (l^q g)^2 \cdot l = (l^q g)^2 \sum_{i=1}^N x_i$. Applying Corollary~\ref{cor-PPconj}, we have
\[ \partial_I W_\Delta(l^p g^2 x_J) = \sum_{i=1}^N  (W_\Delta(l^q g \sqrt{ x_i x_I x_J} ))^2 \]
if the square root exists. Let $L$ have $n-q-m$ components, and let $j$ be one of these components of $L$. Choose $I$ and $J$ such that $x_I x_J = x_L^2/x_j$. Then $\sqrt{ x_j x_I x_J}  = x_L$ and $\sqrt{ x_i x_I x_J}$ does not exist when $i \neq j$, so in this case
\[ \partial_I W_\Delta(l^p g^2 x_J)
 =   (W_\Delta(l^q g x_L))^2 = 0. \]
Hence,
\[ W_\Delta(l^q g x_L) = 0 \]
for any $x_L$. We conclude again that $l^q g = 0$ in $\bar{H}^{q+m}(\Delta)$.
\end{proof}

We now prove a generalization of Theorem~\ref{thm-lef} in the case of pseudo-manifolds.

\begin{theorem}
Let $\Delta$ be a pseudo-manifold of dimension $n-1$ over $K$. Then the quadratic form $Q_l$ on $\bar{H}^m(\Delta)$ is anisotropic  for any $m\leq n/2$.
\end{theorem}

\begin{proof}
Let $g \in \bar{H}^m(\Delta)$ be an isotropic element for $Q_l$,
\[ l^{n-2m} g^2 = 0.\]
Corollary~\ref{cor-reduce} implies that $l^q g = 0$, where $q = \lfloor \frac{n-2m}{2}\rfloor$. In particular, $l^q g^2 = 0$. Continuing this way we reduce the power of $l$ to zero and hence $g=0$.
\end{proof}



The rest of this section consists of the proof of Theorem~\ref{thm-PPconj}.

\subsection{Reductions} We start by reducing Theorem~\ref{thm-PPconj} to simpler cases. First notice that all expressions in Theorem~\ref{thm-PPconj} are defined over the field $\FF_2$. Hence we may assume that $k=\FF_2$. 

Recall that we wrote $W_\Delta = \sum_i W_{\Pi_i}$ in Theorem~\ref{thm-W-sum}. 

\begin{lemma} \label{lem-redPi}
Theorem~\ref{thm-PPconj} for all $\Pi_i$ implies it for $\Delta$.
\end{lemma}

\begin{proof}
This follows directly from the statement of the theorem using the characteristic $2$ assumption and the  observation that if a monomial $x_Ix_J$ restricts to a nonzero monomial on $\Pi_i$, then the monomial is a square if and only if its restriction is a square.
\end{proof}

From now on we will assume that $\Delta = \Pi$ as in Section~\ref{sec-Pi}. Assume that $\Pi$ has vertices $v_0, v_1, \ldots, v_n$. The matrix $A=(\a{i, j})$ has size $n\times (n+1)$, with columns indexed by $0,1,\ldots,n$ and rows by $1,\ldots,n$. We use the notation $X_j$, $j=0,1,\ldots,n$, for the determinant of $A$ with its $j$-th column removed. If $J= (j_1, \ldots, j_s)$ is a vector with entries in $\{0,1,\ldots,n\}$, we let
\[ X_J = X_{j_1} X_{j_2} \cdots X_{j_s}.\]

Theorem~\ref{thm-PPconj} for $\Pi$ can be further reduced to the following:

\begin{theorem} \label{thm-simplest}
Let $I$ and $J$ be vectors with entries in $\{0,1,\dots,n\}$. Assume that $|I|=n$ and $|J|$ is odd.  Then
  \[ \partial_I  X_J  = 
  \begin{cases}
    (X_L)^2 & \text{if $X_I X_J = X_L^2 X_{(0,1,\dots,n)}$}, \\
    0         &  \text{otherwise.}
  \end{cases}\]
\end{theorem}

\begin{lemma} \label{lem-conj-Pi}
Theorem~\ref{thm-simplest} implies Theorem~\ref{thm-PPconj} for $\Pi$.
\end{lemma}
\begin{proof}
Let $I$ and $J$ be as in the statement of Theorem~\ref{thm-PPconj}, and let $J'$ be such that $X_{J'} = X_J X_{(0,1,\dots,n)}$. Note that $|J'|=2n+1$ is odd. We claim that Theorem~\ref{thm-PPconj} for $I,J$ is equivalent to Theorem~\ref{thm-simplest} for $I, J'$. 

Using that 
\[ W_\Pi (x_{J}) = \frac{X_{J}}{c_\Pi}  = \frac{X_{J}}{ X_{(0,1,\dots,n)}}, \]
the statement of Theorem~\ref{thm-PPconj} for $\Pi, I,J$ is
  \[ \partial_I  \frac{X_J}{c_\Pi}  = 
  \begin{cases}
    \frac{X_I X_J}{c_\Pi^2} & \text{if $\sqrt{x_I x_J}$ exists}, \\
    0         &  \text{otherwise.}
  \end{cases}\]
The statement of Theorem~\ref{thm-simplest} for $I, J'$ is 
  \[ \partial_I  ( X_J c_\Pi )  = 
  \begin{cases}
    X_I X_J & \text{if $\sqrt{X_I X_J}$ exists}, \\
    0         &  \text{otherwise.}
  \end{cases}\]
These two equations differ by a factor of $c_\Pi^2$.
\end{proof}

We will prove Theorem~\ref{thm-simplest} below after some preparations.


\subsection{$\SL(n,k)$-invariance}

Let $A=(\a{i, j})$ be the $n\times (n+1)$ matrix of variables. For a matrix $B\in \SL(n,k)$, consider the linear change of variables from $A$ to $BA$. This defines an action of $\SL(n,k)$ on the polynomial ring $k[\a{i, j}]$. The first fundamental theorem of invariant theory for $\SL(n,k)$ states that if $k$ is an infinite field of any characteristic,  then the $k$-algebra of invariants under this action is generated by $X_0, X_1, \dots, X_n$. When  the field $k$ is finite, the same result holds if we consider absolute invariants. These are polynomials in $k[\a{i,j}]$ that are invariant under the action of $\SL(n,\bar{k})$, where $\bar{k}$ is the algebraic closure of $k$. The first fundamental theorem states that absolute invariants are again polynomials in $X_0, X_1, \dots, X_n$ with coefficients in $k$. (See \cite{Procesi2007}, Theorem~13.5.5 and the discussion of absolute invariants in  Section 13.6.1.)

\begin{lemma} \label{lem-invariance}
Let $I$ and $J$ be as in Theorem~\ref{thm-simplest}. Then the polynomial $\partial_I X_J \in k[\a{i, j}]$ is $\SL(n,k)$-invariant for any field $k$ of characteristic $2$. In particular, $\partial_I X_J \in \FF_2[\a{i, j}]$ is an absolute $\SL(n, \FF_2)$-invariant.
\end{lemma}

\begin{proof}
The group $\SL(n,k)$ is generated by elementary matrices. An elementary matrix acts on the matrix of variables by adding a constant $c$ times row $r$ to row $s$. It suffices to prove invariance under this change of variables. 

We may assume without loss of generality that $r=2$ and $s=1$. Consider the new variables
\[ a'_{i,j} = \begin{cases} 
\a{i, j} + c \a{2, j} & \text{if $i=1$,} \\ 
\a{i, j} & \text{otherwise.}
\end{cases}
\]
For a polynomial $f(a) = f(\a{i, j})$,  let us denote by $f(a')$ the result of substituting $\ap{i,j}$ in $\a{i, j}$. Similarly, let us write $\partial_I(a')$ for the partial derivative where we replace $\partial_{\a{i, j}}$ with $\partial_{a'_{i,j}}$. We need to prove that 
\[ (\partial_I X_J)(a') = (\partial_I X_J) (a). \]
We claim that if $\partial_I = \partial_{\a{1, i_1}}  \partial_{\a{2, i_2}} \partial_{\a{3, i_3}}\cdots \partial_{\a{n, i_n}}$, then 
\begin{equation} (\partial_I X_J)(a') = (\partial_I X_J)(a) - c (\partial_{\a{1, i_1}}  \partial_{\a{1, i_2}} \partial_{\a{3, i_3}} \cdots \partial_{\a{n, i_n}} X_J)(a). \label{eq-chain} \end{equation}
The next lemma shows that $\partial_{\a{1, i_1}}  \partial_{\a{1, i_2}} X_J = 0$, hence the second summand vanishes.

For any polynomial $f(\a{i, j})$, by replacing all symbols $a$ with $a'$, we have  
\[ (\partial_I f)(a')= \partial_I(a') f(a').\]
If $f$ is $\SL(n,k)$-invariant, then 
\[ f(a') = f(a) = f(a'_{1,j} - c a'_{2,j}, a'_{2,j},\ldots, a'_{n,j}).\] 
Let us now prove Equation~\eqref{eq-chain}. Since $X_J$ is $\SL(n,k)$-invariant, 
\[  (\partial_I X_J)(a') = \partial_I(a') X_J(a') = \partial_{\ap{1, i_1}}  \partial_{\ap{2, i_2}} \cdots \partial_{\ap{n, i_n}} X_J (a'_{1,j} - c a'_{2,j}, a'_{2,j},\ldots, a'_{n,j}).\]
Using the chain rule, this derivative is
\[ \big( \partial_I X_J - c \partial_{\a{1, i_1}}  \partial_{\a{1, i_2}}\partial_{\a{3, i_3}} \cdots \partial_{\a{n, i_n}} X_J \big) (a'_{1,j} - c a'_{2,j}, a'_{2,j},\ldots, a'_{n,j}).\]
Changing back to the variables $\a{i,j}$ gives the right hand side of  \eqref{eq-chain}.
\end{proof}

\begin{lemma} \label{lem-2deriv}
If $|J|$ is odd then $\partial_{\a{r, i_1}}  \partial_{\a{r, i_2}} X_J = 0$  for any $r, i_1, i_2$.
\end{lemma}

\begin{proof}
  It is enough to consider the case where $J$ contains no repeating indices, since any square factors of $X_J$ can be factored out of the partial derivatives. Under a suitable relabelling of rows and columns of $A$, we can assume that $r=1$, $i_1 = 1$ and $i_2 = 2$, so that the derivative under consideration is $\partial_{\a{1, 1}} \partial_{\a{1, 2}}  X_J$.
  
 Let us denote by $Y_{i,j} = Y_{j,i}$ the determinant of the matrix $A$ with its first row and columns $i,j$ removed. Then
 \[ \partial_{\a{1, i}} X_j = Y_{i,j}.\]
 The polynomials $X_i$ and $Y_{i,j}$ satisfy the following relations. For any distinct indices $i,j,p,q$ in increasing order
 \begin{equation} \label{eq-Plucker}
 Y_{i,j} Y_{p,q} - Y_{i,p} Y_{j,q} + Y_{i,q} Y_{j,p} = 0,
 \end{equation}
 and for any distinct indices $i,j,p$ in increasing order
 \begin{equation} \label{eq-flag}
 Y_{i,j} X_p - Y_{i,p} X_{j} + Y_{j,p} X_{i} = 0.
 \end{equation}
 These equations hold in any characteristic. In characteristic $2$ the signs in the equations are not important.
 The equations come from the Pl\"ucker embedding of the Grassmannian. Rows $2,3,\ldots,n$ of the  matrix $A$ span an $(n-1)$-plane in the $(n+1)$-space and hence define a point in the Grassmannian $\Gr(n-1, n+1)$. The polynomials $Y_{i,j}$ are the Pl\"ucker coordinates on this Grassmannian. These coordinates satisfy the Pl\"ucker relations in Equation~\eqref{eq-Plucker}. Similarly, the $n$ rows of the matrix $A$ define a point in $\Gr(n,n+1)$ with coordinates $X_i$. The relations in Equation~\eqref{eq-flag}  state that the $(n-1)$-plane with coordinates $Y_{i,j}$ lies in the $n$-plane with coordinates $X_i$.
 
Using the product rule we have
\[ \partial_{\a{1, 1}}  \partial_{\a{1, 2}} X_J = \sum \frac{X_J}{X_{i}X_{j}} Y_{1,i} Y_{2,j}.\]
Here the sum runs over all vectors $(i,j)$ where $i,j$ are distinct entries of $J$ such that $i\neq 1$ and $j\neq 2$. We claim that this sum is equal to 
\[ \sum \frac{X_J}{X_{i}X_{j}} Y_{1,2} Y_{i,j},\]
  where the sum now runs over all two element subsets $\{i,j\}$ of entries in $J$. To see this, first consider the case where $i$ and $j$ are both distinct from $1$ and $2$. In this case we apply the Pl\"ucker relation to get 
 \[  Y_{1,i} Y_{2,j} + Y_{2,i} Y_{1,j}= Y_{1,2} Y_{i,j}.\]
 The cases where $i=2$ or $j=1$ are simpler and do not require any relation.
 
 We are now reduced to proving that 
 \[  \sum_{\{i,j\}} \frac{X_J}{X_{i}X_{j}} Y_{1,2} Y_{i,j} = X_J Y_{1,2} \sum_{\{i,j\}} \frac{Y_{i,j}}{X_{i}X_{j}}  = 0.\]
Using Equation~\eqref{eq-flag} we have for any distinct $i,j, p$
\[ \frac{Y_{i,j}}{X_{i}X_{j}} + \frac{Y_{i,p}}{X_{i}X_{p}} +\frac{Y_{j,p}}{X_{j}X_{p}} = 0.\]
Now consider all three element subsets $\{i,j,p\}$ of entries in $J$. Then 
\[ \sum_{\{i,j,p\}} \left( \frac{Y_{i,j}}{X_{i}X_{j}} + \frac{Y_{i,p}}{X_{i}X_{p}} +\frac{Y_{j,p}}{X_{j}X_{p}} \right) = 0.\]
Since every pair $\{i,j\}$ occurs in an odd number of triples $\{i,j,p\}$, this sum is equal to 
\[ \sum_{\{i,j\}} \frac{Y_{i,j}}{X_{i}X_{j}}. \qedhere\]
\end{proof}

\subsection{Proof of Theorem~\ref{thm-simplest}}

Lemma~\ref{lem-invariance} implies that $\partial_I X_J$ is a polynomial in $X_0, X_1, \cdots, X_n$ with coefficients in $\FF_2$.  Consider the grading by $\ZZ^{n+1}$ on the  ring $\FF_2[\a{i, j}]$ such that $\a{i, j}$ has degree $e_j$. Here $e_0,\ldots, e_{n}$ is the standard basis for $\ZZ^{n+1}$. Let $\one = (1,\ldots, 1)$. Then $X_i$, $i=0,\ldots, n$, is homogeneous with
\[ \deg X_i = \one - e_i.\]
Since the vectors $\one - e_i$ are linearly independent, there can be at most one monomial $X_L$ in each degree. 
The partial derivative $\partial_{\a{i,j}}$ applied to a homogeneous polynomial reduces its degree by $e_j$ (or is zero). Since $X_J$ is homogeneous, so is $\partial_I X_J$, hence $\partial_I X_J$ is equal to a constant $c$ times a monomial $X_M$. Here $c\in \FF_2$, hence $\partial_I X_J$ is either $X_M$ or $0$.
Computing the degrees, the monomial $X_M$ must satisfy $X_I X_J = X_M X_{(0,1,\dots,n)}$.
Theorem~\ref{thm-simplest} has two cases depending on whether $\sqrt{X_M}$ exists or not.

\begin{lemma} 
% If $\sqrt{X_M}$ does not exist then $\partial_I X_J = 0$.
If $\partial_I X_J \neq 0$, then $\sqrt{X_M}$ exists.
\end{lemma}

\begin{proof}
Suppose that $\partial_I X_J = X_M$. By Lemma~\ref{lem-2deriv}, all partial derivatives  
 \[ \partial_{\a{i, j}} X_M = \partial_{\a{i, j}} \partial_I X_J\]
 vanish. This implies that $X_M$ is the square of a polynomial in $\FF_2[\a{i, j}]$. Since $X_M$ is the product of irreducible polynomials $X_i$, it follows that $X_M$ must be the square of a monomial $X_L$.
\end{proof}

\begin{lemma}
% If $\sqrt{X_M} = X_L$ then $\partial_I X_J = X_M$.
If $\sqrt{X_M}$ exists, then $\partial_I X_J \neq 0$.
\end{lemma}

\begin{proof}
We will prove that $\partial_I X_J \neq 0$ by induction on $n$. 
The base of the induction is $n=0$. In this case $\partial_I X_J = X_J =1$.

Consider now $n>0$.  Let $I=(i_1,\ldots, i_n)$ and $J=(j_1,\ldots,j_s)$, where $s$ is odd. By assumption, there exists a monomial $X_L$ such that $X_I X_J = X_L^2 X_{(0,1,\ldots,n)}$.  To prove that $\partial_I X_J \neq 0$, we may factor out squares in $X_J$ and assume that $J$ has no repeated entries.

There exists an entry in $J$, say $j_1$, such that $j_1 \neq i_r$ for $r=1,\ldots, n$. This follows from the fact that $X_{(0,1,\ldots,n)}$ of degree $n+1$ divides $X_I X_J$, but $X_I$ has degree $n$. Define the monomial  
  \[ \mu = \a{1,i_1} \a{1,j_1}^{s-1}.\]
Then the coefficient of $\mu$ in $X_J=X_{j_1} \cdots X_{j_s}$ is 
\[  Y_{j_1, i_1} Y_{j_2, j_1} \cdots Y_{j_s,j_1}.\]
Indeed, $X_{j_1}$ does not contain $\a{1,j_1}$. Hence $\a{1,j_1}^{s-1}$ must come from the factors $X_{j_2}, \ldots, X_{j_s}$  and $\a{1,i_1}$ from the factor $X_{j_1}$. As before, we have denoted by $Y_{i,j}$ the coefficient of $\a{1,j}$ in $X_i$.

Notice that the coefficient of $\a{1,j_1}^{s-1}$ in $\partial_I X_J$ is 
\[ \partial_{\a{2,i_2}} \partial_{\a{3,i_3}} \cdots \partial_{\a{n,i_n}} Y_{j_1, i_1} Y_{j_2, j_1} \cdots Y_{j_s,j_1}.\]
If this coefficient is nonzero then also $\partial_I X_J$ is nonzero. We claim that this coefficient being nonzero follows by induction from the case of dimension $n-1$.  From the matrix $(\a{i,j})$ we have removed row $1$ and column $j_1$, the derivative $\partial_I$ is replaced with $\partial_{I'}$, where $I' = (i_2,\ldots, i_n)$, and $X_J$ is replaced (using a similar notation in dimension $n-1$) with $X_{J'}$, where $J'=(i_1, j_2, \ldots, j_s)$. 

Let us check that we can apply the induction assumption to prove that $\partial_{I'} X_{J'} \neq 0$. The vectors $I'$ and $J'$ satisfy
\[ X_{I'} X_{J'} = \frac{X_I}{X_{i_1}} \cdot \frac{ X_{i_1} X_J }{ X_{j_1} } = \frac{X_L^2 X_{(0,1,\dots,n-1,n)}}{X_{j_1}}  \]
Since we assumed that $J$ does not contain repeated entries, $X_{j_1}$ appears in the numerator of the fraction to the first power. If we suppose that $j_1 = n$, then
\[ X_{I'} X_{J'} = X_L^2 X_{(0,1,\dots, n-1)}, \]
where $X_{j_1} = X_n$ does not appear in any monomial. By induction, $\partial_{I'} X_{J'} \neq 0$.
\end{proof}
